\documentclass{article}
\usepackage{notes}

\setcounter{section}{0}
\renewcommand{\thesection}{5.\arabic{section}}

\begin{document}
\stdtitle{ANALYSIS II}{Exercise Sheet 5}{Jakob Schildhauer}

\section{Derivatives}
\begin{enumerate}[label=(\alph*)]
    \item $x = \pm \frac{\sqrt{2}}{2}, y = \pm 1$
    
    \item \dots 
\end{enumerate}

\section{Lagrange Multipliers}
\begin{enumerate}
    \item $\overline{U}$ is compact (closed and bounded) in $R^3$, therefore $f$ attains extremas. 
    
    \item $\nabla f = 0 \implies x = 2, y = 0$, which is not part of the interior of $U$. 
    Therefore, neither $\sup$ nor $\inf$ are attained. 

    \item $f$ is continuous and $\p U$ is compact, therefore a minimum and maximum exist.
    
    \item Let $g(x, y, z) := x^2 + y^2 + z^2 - 1 \fa x, y, z$. 
    Then $\nabla f = \lambda \nabla g$ gives $x = 1, y = z = 0$ or $x = -1, y = z = 0$, where substituting shows that the first is a maximum and the second a minimum.
    
    \item On $\overline{U}$, the maximum and minimum are as described in \titlefont{4.}. 
\end{enumerate}

\section{Multiple Choice}
Correct: \titlefont{(a)}, \titlefont{(c)}.

\section{Computation of Derivatives}
\begin{enumerate}
    \item \begin{itemize}
        \item $\p_v \frac{x_1}{x_2} = v \cdot \left(\frac{1}{x_2}, \frac{-x_1}{x_2^2} \right)^T$
        \item $\p_v e^{-x_1 / x_2} = v \cdot \left( \right)^T$
        \item $\p_v f = v \cdot \left( \right)^T$
        \item $\p_v f = v \cdot \left( \right)^T$
    \end{itemize}
    \item \begin{itemize}
        \item $\p_v f = $
        \item $\p_v f = $
        \item $\p_v f = $
        \item $\p_v f = $
    \end{itemize}
    \item \begin{itemize}
        \item $\p_v f = $
        \item $\p_v f = $
        \item $\p_v f = $
        \item $\p_v f = $
    \end{itemize}
    \item \begin{itemize}
        \item $\p_v f = $
        \item $\p_v f = $
        \item $\p_v f = $
        \item $\p_v f = $
    \end{itemize}
\end{enumerate}

\section{Multiple Choice}
Correct: \titlefont{(a)}, \titlefont{(b)}. 

\section{Polynomials}
Assume by contradiction $\ex \alpha < \sigma$ with $a_{\alpha} \neq b_{\alpha}$ and define $c_k := a_k - b_k$ for all $k \in \N_0$. Then $c_{\alpha} \neq 0$. We have 

\section{Chain Rule for a Curve}
\begin{enumerate}
    \item After some pointless calculations \dots \\
    $r''(t) = R\begin{pmatrix} -\Bigl(b''(t)\sin b(t) + \bigl(b'(t)^2+(\omega+a'(t))^2\bigr)\cos b(t)\Bigr)\cos(\omega t+a(t)) - a''(t)\cos b(t)\sin(\omega t+a(t)) \\[6pt] -\Bigl(b''(t)\sin b(t) + \bigl(b'(t)^2+(\omega+a'(t))^2\bigr)\cos b(t)\Bigr)\sin(\omega t+a(t)) + a''(t)\cos b(t)\cos(\omega t+a(t)) \\[6pt] b''(t)\cos b(t) - b'(t)^2\sin b(t) \end{pmatrix}$
    
    \item \dots
\end{enumerate}
\end{document}